\documentclass[a4paper,12pt]{article}

% 页面设置
\usepackage[top=2.54cm, bottom=2.54cm, left=1.91cm, right=1.91cm]{geometry}

% 中文支持
\usepackage[UTF8]{ctex}
\usepackage{fontspec}

% 数学公式支持
\usepackage{amsmath}
\usepackage{amssymb}
\usepackage{amsfonts}
\usepackage{amsthm}

\usepackage{enumitem}

% 定理环境定义
\newtheorem*{pf}{证}
\newtheorem*{sol}{解}

% 图形支持
\usepackage{graphicx}
\usepackage{tikz}

% 超链接支持
\usepackage{hyperref}
\hypersetup{
    colorlinks=true,
    linkcolor=blue,
    filecolor=magenta,
    urlcolor=cyan,
}

% 行距设置
\linespread{1.25}

% 标题信息
\title{Mathematical Analysis II\\Week 5 Homework (March 31)}
\author{袁子强\hspace{1cm} 2025533009}
% \date{}

\begin{document}

\maketitle

\section*{Section 9.2}

29. 若 $u = F(x,y)$，$F$ 任意二阶偏导存在，而 $x = r\cos\varphi,\,y = r\sin\varphi$。证明：
$$
    \left(\frac{\partial u}{\partial r}\right)^2 + \left(\frac{1}{r}\frac{\partial u}{\partial \varphi}\right)^2 = \left(\frac{\partial u}{\partial x}\right)^2 + \left(\frac{\partial u}{\partial y}\right)^2.
$$

\begin{pf}
    由链式法则得
    $$
        \frac{\partial u}{\partial r}=\frac{\partial u}{\partial x}\frac{\partial x}{\partial r}+\frac{\partial u}{\partial y}\frac{\partial y}{\partial r}=F_x'\cos\varphi+F_y'\sin\varphi,
    $$
    $$
        \frac{\partial u}{\partial \varphi}=\frac{\partial u}{\partial x}\frac{\partial x}{\partial \varphi}+\frac{\partial u}{\partial y}\frac{\partial y}{\partial \varphi}=-F_x'r\sin\varphi+F_y'r\cos\varphi.
    $$
    因此
    $$
        \frac{1}{r}\frac{\partial u}{\partial \varphi}=-F_x'\sin\varphi+F_y'\cos\varphi.
    $$
    从而
    $$
        \begin{aligned}
            \text{LHS} & =F_x'^2\cos^2\varphi+2F_x'F_y'\cos\varphi\sin\varphi+F_y'^2\sin^2\varphi       \\
                       & \quad +F_x'^2\sin^2\varphi-2F_x'F_y'\sin\varphi\cos\varphi+F_y'^2\cos^2\varphi \\
                       & =(F_x'^2+F_y'^2)(\cos^2\varphi+\sin^2\varphi)                                  \\
                       & =F_x'^2+F_y'^2=\text{RHS}.
        \end{aligned}
    $$

    Q.E.D.
\end{pf}

30. 试证：方程 $\dfrac{\partial^2 u}{\partial x^2} + 2\dfrac{\partial^2 u}{\partial x\partial y} - 3\dfrac{\partial^2 u}{\partial y^2} + 2\dfrac{\partial u}{\partial x} + 6\dfrac{\partial u}{\partial y} = 0$ 经变化 $\xi = x + y,\,\eta = 3x - y$ 后变成 $\dfrac{\partial^2 u}{\partial \eta\partial \xi} + \dfrac{1}{2}\dfrac{\partial u}{\partial \xi} = 0$。（其中二阶偏导数均连续）

\begin{pf}
    计算偏导数：
    $$
        \frac{\partial\xi}{\partial x}=1,\quad \frac{\partial\xi}{\partial y}=1,\quad \frac{\partial\eta}{\partial x}=3,\quad \frac{\partial\eta}{\partial y}=-1.
    $$
    因此
    $$
        \frac{\partial u}{\partial x}=u_\xi+3u_\eta,\quad \frac{\partial u}{\partial y}=u_\xi-u_\eta.
    $$
    进而求得二阶偏导数：
    $$
        \begin{aligned}
            \frac{\partial^2u}{\partial x^2}         & =u_{\xi\xi}+3u_{\xi\eta}+3u_{\eta\xi}+9u_{\eta\eta}=u_{\xi\xi}+6u_{\xi\eta}+9u_{\eta\eta}, \\
            \frac{\partial^2u}{\partial x\partial y} & =u_{\xi\xi}+3u_{\xi\eta}-u_{\eta\xi}-3u_{\eta\eta}=u_{\xi\xi}+2u_{\xi\eta}-3u_{\eta\eta},  \\
            \frac{\partial^2u}{\partial y^2}         & =u_{\xi\xi}-u_{\xi\eta}-u_{\eta\xi}+u_{\eta\eta}=u_{\xi\xi}-2u_{\xi\eta}+u_{\eta\eta}.
        \end{aligned}
    $$
    代入原方程，得
    $$
        \begin{aligned}
            \text{LHS} & =u_{\xi\xi}+6u_{\xi\eta}+9u_{\eta\eta}+2(u_{\xi\xi}+2u_{\xi\eta}-3u_{\eta\eta})-3(u_{\xi\xi}-2u_{\xi\eta}+u_{\eta\eta}) \\
                       & \quad +2(u_\xi+3u_\eta)+6(u_\xi-u_\eta)                                                                                 \\
                       & =16u_{\xi\eta}+8u_{\xi}.
        \end{aligned}
    $$
    即
    $$
        16\dfrac{\partial^2 u}{\partial \eta\partial \xi} + 8\dfrac{\partial u}{\partial \xi} = 0.
    $$
    两边同时除以 $16$，即得证。

    Q.E.D.
\end{pf}

32. 设变换 $\begin{cases} u = x - 2y, \\ v = x + ay \end{cases}$ 可把方程 $6\dfrac{\partial^2 z}{\partial x^2} + \dfrac{\partial^2 z}{\partial x\partial y} - \dfrac{\partial^2 z}{\partial y^2} = 0$ 简化为 $\dfrac{\partial^2 z}{\partial u\partial v} = 0$。求常数 $a$。（其中二阶偏导数均连续）

\begin{sol}
    计算偏导数：
    $$
        \frac{\partial u}{\partial x}=1,\quad \frac{\partial u}{\partial y}=-2,\quad \frac{\partial v}{\partial x}=1,\quad \frac{\partial v}{\partial y}=a.
    $$
    因此
    $$
        \frac{\partial z}{\partial x}=z_u+z_v,\quad \frac{\partial z}{\partial y}=-2z_u+az_v.
    $$
    进而求得二阶偏导数：
    $$
        \begin{aligned}
            \frac{\partial^2z}{\partial x^2}         & =z_{uu}+2z_{uv}+z_{vv},        \\
            \frac{\partial^2z}{\partial x\partial y} & =-2z_{uu}+(a-2)z_{uv}+az_{vv}, \\
            \frac{\partial^2z}{\partial y^2}         & =4z_{uu}-4az_{uv}+a^2z_{vv}.
        \end{aligned}
    $$
    代入原方程，得
    $$
        \begin{aligned}
            \text{LHS} & =6(z_{uu}+2z_{uv}+z_{vv})+(-2z_{uu}+(a-2)z_{uv}+az_{vv})-(4z_{uu}-4az_{uv}+a^2z_{vv}) \\
                       & =(10+5a)z_{uv}+(6+a-a^2)z_{vv}.
        \end{aligned}
    $$
    为使方程简化为 $\dfrac{\partial^2 z}{\partial u\partial v} = 0$，需满足
    $$
        \begin{cases}
            6+a-a^2=0, \\
            10+5a\neq 0.
        \end{cases}
    $$
    解得 $a=3$。
\end{sol}

33. 求方程 $\dfrac{\partial z}{\partial y} = x^2 + 2y$ 满足条件 $z(x,x^2) = 1$ 的解 $z = z(x,y)$。

\begin{sol}
    对 $y$ 积分得
    $$
        z = x^2y + y^2 + C(x).
    $$
    代入边界条件 $z(x,x^2) = 1$，得
    $$
        x^2 \cdot x^2 + (x^2)^2 + C(x) = 1,
    $$
    即
    $$
        x^4 + x^4 + C(x) = 1.
    $$
    因此
    $$
        C(x) = 1 - 2x^4.
    $$
    故所求解为
    $$
        z(x,y) = x^2y + y^2 - 2x^4 + 1.
    $$
\end{sol}

34. 设 $u = u(x,y)$，当 $y = x^2$ 时有 $u = 1,\,\dfrac{\partial u}{\partial x} = x$，求当 $y = x^2$ 时的 $\dfrac{\partial u}{\partial y}$。

\begin{sol}
    构造辅助函数 $\varphi(x) = u(x,x^2) = 1$，则 $\varphi'(x) = 0$。

    由链式法则得
    $$
        \begin{aligned}
            \varphi'(x) & = \left.\frac{\partial u}{\partial x}\right|_{y=x^2} + \left.\frac{\partial u}{\partial y}\right|_{y=x^2} \cdot \frac{dy}{dx} \\
                        & = x + \left.\frac{\partial u}{\partial y}\right|_{y=x^2} \cdot \frac{d}{dx}(x^2)                                              \\
                        & = x + \left.\frac{\partial u}{\partial y}\right|_{y=x^2} \cdot 2x = 0.
        \end{aligned}
    $$
    当 $x \neq 0$ 时，两边除以 $2x$，解得
    $$
        \left.\frac{\partial u}{\partial y}\right|_{y=x^2} = -\frac{1}{2}.
    $$
\end{sol}

35. 设 $u = u(x,y)$ 满足方程 $\dfrac{\partial^2 u}{\partial x^2} - \dfrac{\partial^2 u}{\partial y^2} = 0$ 以及条件 $u(x,2x) = x,\,u'_x(x,2x) = x^2$，求 $u''_{xx}(x,2x),\,u''_{xy}(x,2x),\,u''_{yy}(x,2x)$。（其中二阶偏导数均连续）

\begin{sol}
    由方程知
    $$
        u_{xx} = u_{yy}.
    $$
    构造辅助函数 $\varphi(x) = u(x,2x) = x$，则
    $$
        \varphi'(x) = u_x + 2u_y = x^2 + 2u_y = 1,
    $$
    解得
    $$
        u_y = \frac{1 - x^2}{2}.
    $$
    进而
    $$
        \varphi''(x) = u_{xx} + 2u_{xy} + 2u_{yx} + 4u_{yy} = 5u_{xx} + 4u_{xy} = 0.
    $$
    又令 $\psi(x) = u_x(x,2x) = x^2$，则
    $$
        \psi'(x) = u_{xx} + 2u_{xy} = 2x.
    $$
    联立方程组
    $$
        \begin{cases}
            5u_{xx} + 4u_{xy} = 0, \\
            u_{xx} + 2u_{xy} = 2x,
        \end{cases}
    $$
    解得
    $$
        \begin{cases}
            u_{xx} = -\dfrac{4}{3}x, \\[6pt]
            u_{xy} = \dfrac{5}{3}x.
        \end{cases}
    $$
    因此
    $$
        u''_{xx}(x,2x) = -\frac{4}{3}x,\quad u''_{xy}(x,2x) = \frac{5}{3}x,\quad u''_{yy}(x,2x) = u''_{xx}(x,2x) = -\frac{4}{3}x.
    $$
\end{sol}

36. 求下列复合函数的微分 $\mathrm{d}u$

(2) $u = f(\xi,\eta),\,\xi = xy,\,\eta = \dfrac{x}{y}$;

\begin{sol}
    计算中间变量的微分：
    $$
        \mathrm{d}\xi = y\mathrm{d}x + x\mathrm{d}y, \quad \mathrm{d}\eta = \frac{1}{y}\mathrm{d}x - \frac{x}{y^2}\mathrm{d}y.
    $$
    因此
    $$
        \begin{aligned}
            \mathrm{d}u & = \frac{\partial u}{\partial\xi}\mathrm{d}\xi+\frac{\partial u}{\partial\eta}\mathrm{d}\eta                   \\
                        & = f_\xi(y\mathrm{d}x + x\mathrm{d}y) + f_\eta\left(\frac{1}{y}\mathrm{d}x - \frac{x}{y^2}\mathrm{d}y\right)   \\
                        & = \left(f_\xi y + \frac{f_\eta}{y}\right)\mathrm{d}x + \left(f_\xi x - \frac{f_\eta x}{y^2}\right)\mathrm{d}y \\
                        & = \left(f_\xi y + \frac{f_\eta}{y}\right)\mathrm{d}x + \left(f_\xi - \frac{f_\eta}{y^2}\right)x\mathrm{d}y.
        \end{aligned}
    $$
\end{sol}

(4) $u = f(x,\xi,\eta),\,\xi = x^2 + y^2,\,\eta = x^2 + y^2 + z^2$;

\begin{sol}
    计算中间变量的微分：
    $$
        \mathrm{d}\xi = 2x\mathrm{d}x + 2y\mathrm{d}y, \quad \mathrm{d}\eta = 2x\mathrm{d}x + 2y\mathrm{d}y + 2z\mathrm{d}z.
    $$
    因此
    $$
        \begin{aligned}
            \mathrm{d}u & = f_x\mathrm{d}x + f_\xi\mathrm{d}\xi + f_\eta\mathrm{d}\eta                                                    \\
                        & = f_x\mathrm{d}x + f_\xi(2x\mathrm{d}x + 2y\mathrm{d}y) + f_\eta(2x\mathrm{d}x + 2y\mathrm{d}y + 2z\mathrm{d}z) \\
                        & = (f_x + 2xf_\xi + 2xf_\eta)\mathrm{d}x + 2(f_\xi + f_\eta)y\mathrm{d}y + 2zf_\eta\mathrm{d}z.
        \end{aligned}
    $$
\end{sol}

37. 求在球坐标下的 Laplace 方程的形式。

\begin{sol}
    在直角坐标系下，Laplace 方程为
    $$
        \nabla^2u = \frac{\partial^2u}{\partial x^2} + \frac{\partial^2u}{\partial y^2} + \frac{\partial^2u}{\partial z^2} = 0.
    $$
    
    在球坐标系下，坐标变换关系为
    $$
        \begin{cases}
            r = \sqrt{x^2 + y^2 + z^2}, \\
            \theta = \arctan\dfrac{\sqrt{x^2 + y^2}}{z}, \\
            \varphi = \arctan\dfrac{y}{x}.
        \end{cases}
    $$
    由链式法则，
    $$
        \frac{\partial u}{\partial x} = \frac{\partial u}{\partial r}\frac{\partial r}{\partial x} + \frac{\partial u}{\partial \theta}\frac{\partial \theta}{\partial x} + \frac{\partial u}{\partial \varphi}\frac{\partial \varphi}{\partial x} = \frac{\partial u}{\partial r}\frac{x}{r} + \frac{\partial u}{\partial \theta}\frac{xz}{r^2\sqrt{x^2 + y^2}} - \frac{\partial u}{\partial \varphi}\frac{y}{x^2 + y^2}.
    $$
    同理可得
    $$
        \frac{\partial u}{\partial y} = \frac{\partial u}{\partial r}\frac{y}{r} + \frac{\partial u}{\partial \theta}\frac{yz}{r^2\sqrt{x^2 + y^2}} + \frac{\partial u}{\partial \varphi}\frac{x}{x^2 + y^2},
    $$
    $$
        \frac{\partial u}{\partial z} = \frac{\partial u}{\partial r}\frac{z}{r} - \frac{\partial u}{\partial \theta}\frac{\sqrt{x^2 + y^2}}{r^2}.
    $$
    
    由于
    $$
        \begin{cases}
            x = r\sin\theta\cos\varphi, \\
            y = r\sin\theta\sin\varphi, \\
            z = r\cos\theta,
        \end{cases}
    $$
    代入上述表达式，得
    $$
        \frac{\partial u}{\partial x} = \frac{\partial u}{\partial r}\sin\theta\cos\varphi + \frac{\partial u}{\partial \theta}\frac{\cos\theta\cos\varphi}{r} - \frac{\partial u}{\partial \varphi}\frac{\sin\varphi}{r\sin\theta},
    $$
    $$
        \frac{\partial u}{\partial y} = \frac{\partial u}{\partial r}\sin\theta\sin\varphi + \frac{\partial u}{\partial \theta}\frac{\cos\theta\sin\varphi}{r} + \frac{\partial u}{\partial \varphi}\frac{\cos\varphi}{r\sin\theta},
    $$
    $$
        \frac{\partial u}{\partial z} = \frac{\partial u}{\partial r}\cos\theta - \frac{\partial u}{\partial \theta}\frac{\sin\theta}{r}.
    $$
    
    写成矩阵形式，即
    $$
        \begin{pmatrix}
            \dfrac{\partial u}{\partial x} \\[8pt]
            \dfrac{\partial u}{\partial y} \\[8pt]
            \dfrac{\partial u}{\partial z}
        \end{pmatrix} = \underset{\substack{\downarrow\\A}}{
            \begin{pmatrix}
                \sin\theta\cos\varphi & \dfrac{\cos\theta\cos\varphi}{r} & -\dfrac{\sin\varphi}{r\sin\theta} \\[8pt]
                \sin\theta\sin\varphi & \dfrac{\cos\theta\sin\varphi}{r} & \dfrac{\cos\varphi}{r\sin\theta} \\[8pt]
                \cos\theta & -\dfrac{\sin\theta}{r} & 0
            \end{pmatrix}
        } \begin{pmatrix}
            \dfrac{\partial u}{\partial r} \\[8pt]
            \dfrac{\partial u}{\partial \theta} \\[8pt]
            \dfrac{\partial u}{\partial \varphi}
        \end{pmatrix}.
    $$
    
    记
    $$
        \nabla_{xyz} = \begin{pmatrix}
            \dfrac{\partial}{\partial x} \\[6pt]
            \dfrac{\partial}{\partial y} \\[6pt]
            \dfrac{\partial}{\partial z}
        \end{pmatrix}, \quad 
        \nabla_{r\theta\varphi} = \begin{pmatrix}
            \dfrac{\partial}{\partial r} \\[6pt]
            \dfrac{\partial}{\partial \theta} \\[6pt]
            \dfrac{\partial}{\partial \varphi}
        \end{pmatrix}.
    $$
    因此 $\nabla_{xyz}u = A\nabla_{r\theta\varphi}u$。
    
    由 $\nabla^2u = \nabla_{xyz}^T\nabla_{xyz}u$，得
    $$
        \nabla^2u = \nabla_{r\theta\varphi}^T A^T A \nabla_{r\theta\varphi}u + \sum_{1\leq i,j\leq 3} (\partial_i A_{ij}) \partial_j.
    $$
    
    经过计算，得到
    $$
        A^T A = \begin{pmatrix}
            1 & 0 & 0 \\
            0 & \dfrac{1}{r^2} & 0 \\
            0 & 0 & \dfrac{1}{r^2\sin^2\theta}
        \end{pmatrix}.
    $$
    
    记 $B = \displaystyle\sum_{i=1}^3\sum_{j=1}^3 (\partial_i A_{ij})\partial_j$，则
    $$
        B(u) = \frac{2}{r}\frac{\partial u}{\partial r}.
    $$
    
    因此
    \begin{align*}
        \nabla^2u & = \nabla_{r\theta\varphi}^T A^T A \nabla_{r\theta\varphi}u + B(u) \\
                  & = \begin{pmatrix}
                          \dfrac{\partial}{\partial r} & \dfrac{\partial}{\partial \theta} & \dfrac{\partial}{\partial \varphi}
                      \end{pmatrix}
                      \begin{pmatrix}
                          1 & 0 & 0 \\
                          0 & \dfrac{1}{r^2} & 0 \\
                          0 & 0 & \dfrac{1}{r^2\sin^2\theta}
                      \end{pmatrix}
                      \begin{pmatrix}
                          \dfrac{\partial u}{\partial r} \\[6pt]
                          \dfrac{\partial u}{\partial \theta} \\[6pt]
                          \dfrac{\partial u}{\partial \varphi}
                      \end{pmatrix} + \frac{2}{r}\frac{\partial u}{\partial r} \\
                  & = \frac{\partial^2 u}{\partial r^2} + \frac{1}{r^2}\frac{\partial^2 u}{\partial \theta^2} + \frac{1}{r^2\sin^2\theta}\frac{\partial^2 u}{\partial \varphi^2} + \frac{2}{r}\frac{\partial u}{\partial r}.
    \end{align*}
    
    由于
    $$
        \frac{\partial^2 u}{\partial r^2} + \frac{2}{r}\frac{\partial u}{\partial r} = \frac{1}{r^2}\frac{\partial}{\partial r}\left(r^2\frac{\partial u}{\partial r}\right),
    $$
    因此球面坐标系下的 Laplace 方程的形式为
    $$
        \nabla^2u = \frac{1}{r^2}\frac{\partial}{\partial r}\left(r^2\frac{\partial u}{\partial r}\right) + \frac{1}{r^2}\frac{\partial^2 u}{\partial \theta^2} + \frac{1}{r^2\sin^2\theta}\frac{\partial^2 u}{\partial \varphi^2} = 0.
    $$
\end{sol}

\section*{Section 9.3}

1. 证明下列方程在指定点的附近对 $y$ 有唯一解，并求出 $y$ 对 $x$ 在该点处的一阶和二阶导数

(1) $x^2 + xy + y^2 = 7$，在 $(2,1)$ 处.

\begin{pf}
    令 $F(x,y) = x^2 + xy + y^2 - 7 \in C^1$，则 $F(2,1) = 4 + 2 + 1 - 7 = 0$。

    计算偏导数得 $F_y = x + 2y$，因此 $F_y(2,1) = 4 \neq 0$。由隐函数存在定理知，在 $(2,1)$ 附近对 $y$ 有唯一解。

    Q.E.D.
\end{pf}

\begin{sol}
    由隐函数求导公式得
    $$
        \frac{\mathrm{d}y}{\mathrm{d}x} = -\frac{F_x}{F_y} = -\frac{2x + y}{x + 2y}.
    $$
    因此在点 $(2,1)$ 处，
    $$
        y' = \left.\frac{\mathrm{d}y}{\mathrm{d}x}\right|_{(2,1)} = -\frac{5}{4}.
    $$
    对 $\dfrac{\mathrm{d}y}{\mathrm{d}x}$ 关于 $x$ 求导，得
    $$
        y'' = -\frac{(2 + y')(x + 2y) - (2x + y)(1 + 2y')}{(x + 2y)^2}.
    $$
    代入 $(2,1)$ 及 $y' = -\dfrac{5}{4}$，计算得
    $$
        y''|_{(2,1)} = -\frac{\frac{3}{4} \cdot 4 + 5 \cdot \frac{3}{2}}{4^2} = -\frac{21}{32}.
    $$
\end{sol}

(2) $x\cos xy = 0$，在 $(1,\frac{\pi}{2})$ 处.

\begin{pf}
    令 $F(x,y) = x\cos xy \in C^1$，则 $F(1,\frac{\pi}{2}) = \cos\frac{\pi}{2} = 0$。

    计算偏导数得 $F_y = -x^2\sin xy$，因此 $F_y(1,\frac{\pi}{2}) = -1 \neq 0$。由隐函数存在定理知，在 $(1,\frac{\pi}{2})$ 附近对 $y$ 有唯一解。

    Q.E.D.
\end{pf}

\begin{sol}
    由隐函数求导公式得
    $$
        \frac{\mathrm{d}y}{\mathrm{d}x} = -\frac{F_x}{F_y} = \frac{\cos xy - xy\sin xy}{x^2\sin xy}.
    $$
    因此在点 $(1,\frac{\pi}{2})$ 处，
    $$
        y' = \left.\frac{\mathrm{d}y}{\mathrm{d}x}\right|_{(1,\frac{\pi}{2})} = -\frac{\pi}{2}.
    $$
    对 $\dfrac{\mathrm{d}y}{\mathrm{d}x}$ 关于 $x$ 求导，经过化简得
    $$
        y'' = \frac{2y'}{x} + \frac{-\cos xy(2x\sin xy + x^2(y + xy')\cos xy)}{x^4\sin^2xy}.
    $$
    代入 $(1,\frac{\pi}{2})$ 及 $y' = -\dfrac{\pi}{2}$，计算得
    $$
        y''|_{(1,\frac{\pi}{2})} = -2y' = \pi.
    $$
\end{sol}

2. 求由下列方程所确定的隐函数的导数.

(1) $\sin(xy) - \mathrm{e}^{xy} - x^2y = 0$，求 $\dfrac{\mathrm{d}y}{\mathrm{d}x}$;

\begin{sol}
    方程两边关于 $x$ 求导，得
    $$
        y\cos xy + xy'\cos xy - y\mathrm{e}^{xy} - xy'\mathrm{e}^{xy} - 2xy - x^2y' = 0.
    $$
    整理得
    $$
        xy'(\cos xy - \mathrm{e}^{xy} - x) = -y\cos xy + y\mathrm{e}^{xy} + 2xy.
    $$
    因此
    $$
        \frac{\mathrm{d}y}{\mathrm{d}x} = y' = \frac{-y\cos xy + y\mathrm{e}^{xy} + 2xy}{x\cos xy - x\mathrm{e}^{xy} - x^2}.
    $$
\end{sol}

(3) $x^y = y^x$，求 $\dfrac{\mathrm{d}y}{\mathrm{d}x}$ 和 $\dfrac{\mathrm{d}^2y}{\mathrm{d}x^2}$;

\begin{sol}
    两边取对数得 $y\ln x = x\ln y$。

    两边关于 $x$ 求导，得
    $$
        y'\ln x + \frac{y}{x} = \ln y + \frac{xy'}{y}.
    $$
    整理得
    $$
        y'(y\ln x - x) = x\ln y - y.
    $$
    因此
    $$
        \frac{\mathrm{d}y}{\mathrm{d}x} = y' = \frac{x\ln y - y}{y\ln x - x}.
    $$
    进一步求二阶导数，得
    $$
        \frac{\mathrm{d}^2y}{\mathrm{d}x^2} = \frac{(1 - y')(y'\ln x + \frac{y}{x})}{y\ln x - x}.
    $$
\end{sol}

(5) $\dfrac{x}{z} = \ln\dfrac{z}{y}$，求 $\dfrac{\partial z}{\partial x}$, $\dfrac{\partial z}{\partial y}$;

\begin{sol}
    令 $F(x,y,z) = \dfrac{x}{z} + \ln y - \ln z = 0$。

    计算偏导数得
    $$
        F_x = \frac{1}{z}, \quad F_y = \frac{1}{y}, \quad F_z = -\frac{x}{z^2} - \frac{1}{z}.
    $$
    由隐函数求导公式得
    $$
        \frac{\partial z}{\partial x} = -\frac{F_x}{F_z} = \frac{\frac{1}{z}}{\frac{x}{z^2} + \frac{1}{z}} = \frac{z}{x + z},
    $$
    $$
        \frac{\partial z}{\partial y} = -\frac{F_y}{F_z} = \frac{\frac{1}{y}}{\frac{x}{z^2} + \frac{1}{z}} = \frac{z^2}{(x + z)y}.
    $$
\end{sol}

(7) $F(xz,yz) = 0$，求 $\dfrac{\partial z}{\partial x}$, $\dfrac{\partial z}{\partial y}$.

\begin{sol}
    令 $u = xz$, $v = yz$。

    方程两边关于 $x$ 求偏导，得
    $$
        F_u\left(z + x\frac{\partial z}{\partial x}\right) + F_v\left(y\frac{\partial z}{\partial x}\right) = 0.
    $$
    解得
    $$
        \frac{\partial z}{\partial x} = -\frac{F_u z}{F_u x + F_v y}.
    $$
    方程两边关于 $y$ 求偏导，得
    $$
        F_u\left(x\frac{\partial z}{\partial y}\right) + F_v\left(z + y\frac{\partial z}{\partial y}\right) = 0.
    $$
    解得
    $$
        \frac{\partial z}{\partial y} = -\frac{F_v z}{F_u x + F_v y}.
    $$
\end{sol}

\section*{第九章综合习题}

2. 设 $f(x,y,z) = F(u,v,w)$，其中 $x^2 = vw,\,y^2 = wu,\,z^2 = uv$。求证：
$$
    x\frac{\partial f}{\partial x} + y\frac{\partial f}{\partial y} + z\frac{\partial f}{\partial z} = u\frac{\partial F}{\partial u} + v\frac{\partial F}{\partial v} + w\frac{\partial F}{\partial w}.
$$

\begin{pf}
    由 $x^2 = vw$, $y^2 = wu$, $z^2 = uv$ 得 $yz = u\sqrt{vw} = ux$，因此
    $$
        u = \frac{yz}{x}.
    $$
    同理可得
    $$
        v = \frac{xz}{y}, \quad w = \frac{xy}{z}.
    $$
    由链式法则，
    $$
        \frac{\partial f}{\partial x} = \frac{\partial F}{\partial u}\frac{\partial u}{\partial x} + \frac{\partial F}{\partial v}\frac{\partial v}{\partial x} + \frac{\partial F}{\partial w}\frac{\partial w}{\partial x} = -\frac{yz}{x^2}F_u + \frac{z}{y}F_v + \frac{y}{z}F_w.
    $$
    因此
    $$
        x\frac{\partial f}{\partial x} = -\frac{yz}{x}F_u + \frac{xz}{y}F_v + \frac{xy}{z}F_w = -uF_u + vF_v + wF_w.
    $$
    同理可得
    $$
        y\frac{\partial f}{\partial y} = uF_u - vF_v + wF_w, \quad z\frac{\partial f}{\partial z} = uF_u + vF_v - wF_w.
    $$
    将上述三式相加，得
    $$
        \text{LHS} = 2uF_u - uF_u + 2vF_v - vF_v + 2wF_w - wF_w = uF_u + vF_v + wF_w = \text{RHS}.
    $$

    Q.E.D.
\end{pf}

3. 若函数 $u = f(x,y,z)$ 满足恒等式 $f(tx,ty,tz) = t^k f(x,y,z)\ (t>0)$，则称 $f(x,y,z)$ 为 $k$ 次齐次函数。试证下述关于齐次函数的欧拉定理：可微函数 $f(x,y,z)$ 为 $k$ 次齐次函数的充要条件是：
$$
    x f'_x(x,y,z) + y f'_y(x,y,z) + z f'_z(x,y,z) = k f(x,y,z).
$$

\begin{pf}
    \begin{itemize}
        \item 必要性

              因为 $f(tx,ty,tz) = t^k f(x,y,z)$，两边关于 $t$ 求导，得
              $$
                  xf_x'(tx,ty,tz) + yf_y'(tx,ty,tz) + zf_z'(tx,ty,tz) = kt^{k-1}f(x,y,z).
              $$
              令 $t = 1$，得
              $$
                  xf_x'(x,y,z) + yf_y'(x,y,z) + zf_z'(x,y,z) = kf(x,y,z).
              $$
              必要性得证。

        \item 充分性

              构造辅助函数 $\varphi(t) = \dfrac{f(tx,ty,tz)}{t^k}$。

              对 $t$ 求导，得
              $$
                  \varphi'(t) = \frac{(xf_x'(tx,ty,tz) + yf_y'(tx,ty,tz) + zf_z'(tx,ty,tz))t^k - kf(tx,ty,tz)t^{k-1}}{t^{2k}}.
              $$
              代入条件 $xf_x'(x,y,z) + yf_y'(x,y,z) + zf_z'(x,y,z) = kf(x,y,z)$，得
              $$
                  \varphi'(t) = \frac{kf(tx,ty,tz)t^k \cdot \frac{1}{t} - kf(tx,ty,tz)t^{k-1}}{t^{2k}} = 0.
              $$
              因此 $\varphi(t)$ 恒为常数，故
              $$
                  \varphi(t) = \frac{f(tx,ty,tz)}{t^k} = \varphi(1) = f(x,y,z),
              $$
              即 $f(tx,ty,tz) = t^k f(x,y,z)$。充分性得证。
    \end{itemize}

    Q.E.D.
\end{pf}

\end{document}